\paragraph{Retained sources and atomic cofactors.}
The projection-profile fiber is part of a larger class with an exact
lower-arity normal form. A source is \emph{retained} when its whole input
menu is contained in every output. The first lemma removes a fixed menu
from a cofactor without treating the empty set as a unit.

\begin{lemma}[Erasing a fixed menu]\label{core:lem:winner-erasure}
For $K\geq3$, fix a core word $t(X,v)$ and a nonempty menu $A$, and put
$h(v)=t(A,v)$ and $C=\partial_Xt$.
Either $h$ is constantly $A$, or there is a core operation $g$ on
the remaining variables such that, at every valuation,
\[
 g(v)\subseteq h(v)\cap C(v),\qquad
 g(v)\setminus A=h(v)\setminus A.
\]
The erasure $g$ preserves every nonzero singleton predicate of $h$.
If $A$ is ambiguous and $C$ is nonempty, one may take $g=C$,
simultaneously for all such menus $A$.
\end{lemma}
\begin{proof}
For ambiguous $A$, the ghost comparison gives either $h=C$ or
$h=A\cup C$ at each valuation. Empty residue gives the constant case;
otherwise take $g=C$. The singleton predicates of $h$ are precisely
those obtained by excluding $X$, so
\cref{core:cor:residual-compiler} supplies the profile assertion.

It remains to consider $A=\{q\}$.
For each subword after fixing $X=\{q\}$, write $\psi_s(E)$ for its
$s$-singleton predicate on the $s$-atomic positions of the remaining
inputs. These predicates are monotone. The $q$-predicate is nonzero,
and $\psi_q(\varnothing)=1$ exactly when the subword is constantly
$\{q\}$. Along with the two inclusions in the statement, maintain
\[
 \Phi_s(g)=\psi_s\qquad\text{whenever }\psi_s\text{ is nonzero}.
\]
An ordinary variable is unchanged. An $X$-leaf is the constant case.

At a root owned by $r$, a constant-$q$ child makes the parent constant
when $r=q$; two constant children also do so. At a foreign root discard
a constant-$q$ child and use the other child's erasure. This preserves
all non-$q$ memberships and every nonzero singleton predicate of the
parent. Suppose now both children are nonconstant and have erasures
$g_1,g_2$. Use the four-row table in the proof of
\cref{core:thm:neutral-elimination}, now testing their fixed-$q$
$r$-predicates. At $r=q$ both are nonzero, giving $g_1\star_qg_2$.
A child with zero $r$-predicate never equals $\{r\}$, although its
erasure might. The added union prevents this artificial absorber from
discarding the other operand. For a child with nonzero predicate, the
induction hypothesis preserves its $r$-singleton test exactly. Hence
each row has precisely the parent's non-$q$ memberships and is contained
in its value. Foreign singleton predicates conjoin; at the root owner's
one-zero case the profile is
$(\Phi_r(g_1)\land\Phi_r(g_2))\lor\Phi_r(g_2)=\Phi_r(g_2)$,
or symmetrically. This proves the profile invariant as well.

For containment in $C$, use the unowned ghost from
\cref{core:thm:neutral-elimination}. If the ghost root absorbs at $r$,
some child is unmarked $\{r\}$, so its fixed-$q$ value is also
$\{r\}$. Its nonzero predicate makes the reconstructed erasure absorb
at $r$ too. Otherwise the ordinary residue is the union of the child
residues, which contains the reconstructed erasure by induction. In the
discarded-constant case at $r\ne q$, that child's ghost can only be
marked or unmarked $\{q\}$: an unmarked ghost is unchanged by replacing
it with $\{q\}$. It therefore cannot absorb at $r$. Any absorption by
the other child is preserved; otherwise its residue is contained in the
parent's. This completes the induction.
\end{proof}

\begin{theorem}[Normal form for a retained source]
\label{core:thm:retained-source}
Fix a finite owner set of size $K\geq3$. A core operation $F(X,v)$
satisfying $X\subseteq F(X,v)$ is either the projection $X$, or has
a core operation $C(v)$ and, for each owner $q$, an optional core
operation $g_q(v)\leq C(v)$ such that
\[
 F(X,v)=
 \begin{cases}
  \{q\}\cup g_q(v),&X=\{q\},\\
  X\cup C(v),&X\text{ ambiguous}.
 \end{cases}
\]
An omitted $g_q$ contributes nothing. The baseline is intrinsic:
$C=\partial_XF$. Conversely every such choice of lower-arity core
operations defines a core operation retaining $X$.
Every operation in this class has a primitive core realization with
at most $2K^2$ occurrences of $X$, independently of its arity.
\end{theorem}
\begin{proof}
Replace a realizing word $t$ by $t+X$, which has the same operation.
Its ghost always survives. Neutral-source elimination therefore gives
$F(X,v)=X\cup C(v)$ on ambiguous inputs. Empty residue gives the
projection; otherwise $C$ is core. Apply
\cref{core:lem:winner-erasure} at each atom $\{q\}$.
Retention turns its non-$q$ equality into
$F(\{q\},v)=\{q\}\cup g_q(v)$, and its residue bound gives $g_q\leq C$.
Constant-$q$ cofactors are the omitted cases.

For the converse and the size bound, start with the baseline piece
$P_0=C\del X$. It has a realization by starting at $C$ and attaching
$X$ once at each of the $K$ owners. At atomic $X$ its owner's stage
absorbs to $X$; at ambiguous $X$, at most the first stage can absorb
to an atom, and the next distinct owner restores $C\cup X$.
Thus this realization uses $K$ occurrences of $X$.
For each nonomitted $g_q$, form $P_q$ by starting at $g_q$ and
attaching $X$ once at every owner other than $q$. Then
\[
 P_q(X,v)=
 \begin{cases}
 X,&X=\{s\},\ s\ne q,\\
 X\cup g_q(v),&X=\{q\}\text{ or }X\text{ ambiguous}.
 \end{cases}
\]
There are at least two attachment stages. If the first absorbs the
starting value, the next distinct owner restores its union with $X$;
at another atom $\{s\}$ the $s$-stage instead absorbs to $X$.
These pieces use at most $K+K(K-1)=K^2$ occurrences of $X$.
Their pointwise union is exactly the displayed normal form.

To implement that union without further replication, choose distinct
owners $a,b,c$. Fold all the pieces by $\star_a$, obtaining $W_a$,
and separately by $\star_b$, obtaining $W_b$. For ambiguous $X$, all
pieces are ambiguous and both folds equal $F$. At $X=\{q\}$ every
piece contains $q$, so $W_a=F$ unless $q=a$; then the baseline piece
forces $W_a=\{a\}$. The analogous assertion holds for $W_b$.
Consequently $W_a\star_cW_b=F$ in every case. Each piece is used
twice, proving the primitive bound $2K^2$.
\end{proof}

A retained source is detected by the canonical coefficients
$c_{k,X}=1$ at every observer. Thus retained sources admit uniformly
bounded reuse even though arbitrary
sources do not (\cref{core:thm:robust-tree-size}). The normal form reduces
this class to lower-arity core operations. Retained extensions still need
not glue: \cref{core:cor:prime-obstructions} gives a noncore $F$ for which
every $F+x_j$ is core, and whose pointwise intersection is $F$.
